% MIKU 多障碍物统一规划算法总流程
\begin{tikzpicture}[
    scale=1.0, transform shape,
    stepimp/.style={rectangle, rounded corners=3pt, draw=deeppink, fill=improvedpink, minimum width=5.4cm, minimum height=1.0cm, font=\small, align=center},
    inp/.style={rectangle, draw=deepblue, fill=inputyellow, minimum width=4.0cm, minimum height=0.8cm, font=\footnotesize, align=center},
    outnode/.style={rectangle, rounded corners=3pt, draw=deeppink, fill=outputpink, minimum width=4.0cm, minimum height=0.8cm, font=\small\bfseries, align=center},
    qpnode/.style={rectangle, rounded corners=3pt, draw=deepblue, fill=lightblue, minimum width=4.0cm, minimum height=0.8cm, font=\small\bfseries, align=center},
    arr/.style={-Stealth, thick, deepblue},
    injarr/.style={-Stealth, thick, deeppink},
]

% 输入
\node[inp] (i1) at (-5.8, 6.0) {参考线 $+$ 障碍物集合};
\node[inp] (i2) at (-5.8, 5.0) {预测轨迹 $+$ 车辆状态};

% 步骤（间距 1.7，全部为 MIKU 改进）
\node[stepimp] (s1) at (0, 5.5)  {\textbf{Step 1} 到达时间 $\tau(s)$ 与时变 SL 投影};
\node[stepimp] (s2) at (0, 3.8)  {\textbf{Step 2} 多因子威胁度评估 $\Theta_i$};
\node[stepimp] (s3) at (0, 2.1)  {\textbf{Step 3} 动态安全裕度 $\delta_i$ 与边界 $u_i,v_i$};
\node[stepimp] (s4) at (0, 0.4)  {\textbf{Step 4} 扫描线重叠分组 $\mathcal{G}_1,\ldots,\mathcal{G}_m$};
\node[stepimp] (s5) at (0,-1.3)  {\textbf{Step 5} 组内最大间隙 $g^\star_j=\max_k g_{j,k}$};
\node[stepimp] (s6) at (0,-3.0)  {\textbf{Step 6} 合并全局路径横向边界};
\node[stepimp] (s7) at (0,-4.7)  {\textbf{Step 7} ST 时空可通行走廊 $\Omega$};

% 输出（左侧）与跨接口的速度 QP
\node[outnode] (o1) at (-5.8,-3.0) {路径横向边界};
\node[outnode] (o2) at (-5.8,-4.7) {时空可通行走廊};
\node[qpnode]  (qp) at (-5.8,-6.4) {速度 QP（求解器不变）};

% 输入箭头
\draw[arr] (i1.east) -- (s1.west);
\draw[arr] (i2.east) -- (s1.west);

% 步骤间箭头
\draw[arr] (s1) -- (s2);
\draw[arr] (s2) -- (s3);
\draw[arr] (s3) -- (s4);
\draw[arr] (s4) -- (s5);
\draw[arr] (s5) -- (s6);
\draw[arr] (s6) -- (s7);

% 输出箭头
\draw[arr] (s6.west) -- (o1.east);
\draw[arr] (s7.west) -- (o2.east);

% 走廊注入速度 QP（跨解耦接口）
\draw[injarr] (o2.south) -- node[right, font=\footnotesize, deeppink] {逐时间步纵向上界} (qp.north);

\end{tikzpicture}
